%========================================
% MatinBook - Stage 03: Layout Test
% Test: Page Geometry + Headers/Footers + Headings + TOC
% Purpose: Validate page layout and navigation systems
%========================================

% Add parent directory to search path
\makeatletter
\def\input@path{{../}}
\makeatother

\documentclass{matinbook}

\title{آزمون صفحه‌آرایی و ساختار}
\author{تیم توسعه MatinBook}
\date{تابستان ۱۴۰۵}

\begin{document}

%----------------------------------------
% Title Page
%----------------------------------------
\maketitle

%----------------------------------------
% Table of Contents
%----------------------------------------
\tableofcontents
\clearpage

%----------------------------------------
% Chapter 1: Page Geometry
%----------------------------------------
\chapter{هندسه صفحه}

\section{تنظیمات حاشیه}

این پاراگراف برای بررسی حاشیه‌های صفحه نوشته شده است.
کلاس \lr{MatinBook} از بسته \lr{geometry} برای تنظیم دقیق
ابعاد صفحه استفاده می‌کند. تنظیمات پیش‌فرض شامل موارد زیر است:

\begin{itemize}
    \item \textbf{اندازه کاغذ:} \lr{A4} (۲۱۰ × ۲۹۷ میلی‌متر)
    \item \textbf{حاشیه داخلی:} ۲.۵ سانتی‌متر
    \item \textbf{حاشیه خارجی:} ۲.۵ سانتی‌متر
    \item \textbf{حاشیه بالا:} ۲.۵ سانتی‌متر
    \item \textbf{حاشیه پایین:} ۲.۵ سانتی‌متر
    \item \textbf{فضای صحافی:} ۱ سانتی‌متر (\lr{bindingoffset})
\end{itemize}

\section{تست متن طولانی}

این یک پاراگراف طولانی است که برای تست نحوه شکسته شدن خطوط
و پر شدن صفحه طراحی شده است. هدف ما این است که ببینیم آیا
متن به درستی در صفحه جای می‌گیرد و حاشیه‌ها رعایت می‌شوند
یا خیر. \lr{MatinBook} از تنظیمات پیشرفته \lr{microtype}
برای بهبود ترازبندی و جلوگیری از سرریز خطوط استفاده می‌کند.

فضای اضافی ۱ سانتی‌متری برای صحافی (\lr{bindingoffset})
در نظر گرفته شده است که برای کتاب‌های چاپی ضروری است.
این فضا باعث می‌شود متن در قسمت داخلی صفحه (نزدیک عطف کتاب)
گم نشود و خوانایی کتاب در حالت باز حفظ گردد.

\section{تست فاصله پاراگراف‌ها}

این پاراگراف اول است. فاصله بین پاراگراف‌ها باید مناسب باشد
--- نه آنقدر زیاد که متن از هم گسیخته به نظر برسد،
و نه آنقدر کم که پاراگراف‌ها در هم ادغام شوند.

این پاراگراف دوم است. تنظیمات پیش‌فرض \lr{MatinBook} شامل
\texttt{\textbackslash parindent=1em} برای تورفتگی ابتدای پاراگراف
و \texttt{\textbackslash parskip=0pt plus 1pt} برای فاصله بین
پاراگراف‌ها می‌باشد.

این پاراگراف سوم است. توجه کنید که تورفتگی فقط برای پاراگراف‌های
دوم به بعد اعمال می‌شود و پاراگراف اول هر بخش تورفتگی ندارد.
این یک استاندارد در حروف‌چینی فارسی است.

\section{پانویس‌ها}

این متن دارای یک پانویس است\footnote{این یک پانویس آزمایشی است
که باید در پایین صفحه نمایش داده شود. فاصله پانویس‌ها از متن
و از یکدیگر باید مناسب باشد. خط جداکننده پانویس باید ۳۰٪
عرض متن باشد.}.

و این هم یک پانویس دیگر\footnote{پانویس دوم برای تست فاصله
بین پانویس‌ها. \texttt{\textbackslash footnotesep=8pt}
تنظیم شده است.} برای اطمینان از صحت تنظیمات.

%----------------------------------------
% Chapter 2: Headers and Footers
%----------------------------------------
\chapter{سربرگ و پابرگ}

\section{سربرگ (\lr{Header})}

این صفحه باید سربرگ مناسب با عنوان فصل و شماره صفحه داشته باشد.
در \lr{MatinBook} از بسته \lr{fancyhdr} برای مدیریت سربرگ‌ها
و پابرگ‌ها استفاده می‌شود. تنظیمات به صورت دوطرفه (\lr{twoside})
انجام شده است:

\begin{itemize}
    \item \textbf{صفحات فرد:} نام فصل در سمت راست، شماره صفحه در چپ
    \item \textbf{صفحات زوج:} شماره صفحه در راست، نام بخش در چپ
    \item \textbf{صفحات شروع فصل:} فقط پابرگ (بدون سربرگ)
\end{itemize}

\section{پابرگ (\lr{Footer})}

پابرگ شامل یک خط افقی نازک و عبارت \lr{MatinBook} به همراه
شماره صفحه است. این طراحی ساده اما حرفه‌ای، به خواننده
کمک می‌کند بدون حواس‌پرتی روی محتوا تمرکز کند.

\section{تست چندین صفحه}

برای تست کامل سربرگ و پابرگ، نیاز به چندین صفحه داریم.
این پاراگراف‌های طولانی برای پر کردن صفحات نوشته شده‌اند.

\subsection{زیربخش اول}

محتوای علمی می‌تواند شامل توضیحات مفصل درباره یک موضوع باشد.
برای مثال، در این بخش می‌توانیم درباره اهمیت صفحه‌آرایی
در کتاب‌های فارسی صحبت کنیم. صفحه‌آرایی خوب باعث می‌شود
خواننده تجربه بهتری از مطالعه داشته باشد و مطالب را
با سرعت و دقت بیشتری درک کند.

\subsection{زیربخش دوم}

یکی از چالش‌های صفحه‌آرایی فارسی، ترکیب متن فارسی با
عناصر لاتین مانند کدهای برنامه‌نویسی، فرمول‌های ریاضی
و جداول است. \lr{MatinBook} با استفاده از \lr{xepersian}
و تنظیمات دقیق، این مشکل را به خوبی حل کرده است.

\subsection{زیربخش سوم}

فاصله بین بخش‌ها و زیربخش‌ها باید به گونه‌ای باشد که
سلسله مراتب مطالب به وضوح مشخص شود. عناوین اصلی باید
برجسته‌تر از عناوین فرعی باشند و خواننده بتواند
به راحتی ساختار مطلب را دنبال کند.

%----------------------------------------
% Chapter 3: Heading Styles
%----------------------------------------
\chapter{استایل عناوین}

\section{عنوان فصل (\lr{Chapter})}

عنوان فصل با استایل خاص \lr{MatinBook} شامل یک شماره بزرگ
و کمرنگ در پس‌زمینه، عنوان با فونت درشت، و یک خط رنگی
در پایین است. این طراحی از کتاب‌های فنی مدرن الهام گرفته
شده است.

\section{عنوان بخش (\lr{Section})}

این یک \texttt{\textbackslash section} است. عنوان بخش
با جعبه رنگی شماره و فونت \lr{Large} نمایش داده می‌شود.

\subsection{عنوان زیربخش (\lr{Subsection})}

این یک \texttt{\textbackslash subsection} با فونت \lr{large}
و شماره آبی کمرنگ است.

\subsubsection{عنوان زیرزیربخش (\lr{Subsubsection})}

این یک \texttt{\textbackslash subsubsection} با فونت
معمولی و شماره آبی خیلی کمرنگ است.

\section{تست سلسله مراتب کامل}

\subsection{سطح دوم: زیربخش}

این متن زیر یک زیربخش قرار دارد. باید تورفتگی مناسبی
نسبت به عنوان بخش داشته باشد.

\subsubsection{سطح سوم: زیرزیربخش}

این متن زیر یک زیرزیربخش قرار دارد. عمیق‌ترین سطح
عنوان‌بندی در \lr{MatinBook} است.

\subsection{زیربخش دیگر}

برای مقایسه فاصله بین دو زیربخش متوالی.

%----------------------------------------
% Chapter 4: Table of Contents
%----------------------------------------
\chapter{فهرست مطالب}

\section{تنظیمات فهرست}

فهرست مطالب \lr{MatinBook} دارای ویژگی‌های زیر است:

\begin{itemize}
    \item \textbf{فونت فصل‌ها:} سیاه و پررنگ
    \item \textbf{شماره صفحه فصل‌ها:} آبی پررنگ
    \item \textbf{نقطه‌چین:} با رنگ آبی کمرنگ
    \item \textbf{تورفتگی:} برای سطوح مختلف
    \item \textbf{فاصله:} ۸pt قبل از فصل، ۴pt قبل از بخش
\end{itemize}

\section{فهرست تصاویر و جداول}

علاوه بر فهرست مطالب، \lr{MatinBook} از فهرست تصاویر
(\texttt{\textbackslash listoffigures})، فهرست جداول
(\texttt{\textbackslash listoftables}) و فهرست الگوریتم‌ها
(\texttt{\textbackslash listofalgorithms}) نیز پشتیبانی می‌کند.

\section{نمونه ارجاع به فهرست}

می‌توانیم به بخش‌های مختلف ارجاع دهیم. برای مثال،
در \chref{chap:summary} نتایج نهایی را مشاهده خواهید کرد.

%----------------------------------------
% Chapter 5: Page Styles
%----------------------------------------
\chapter{استایل‌های صفحه}

\section{صفحه شروع فصل}

صفحات شروع فصل (\lr{plain style}) فاقد سربرگ هستند
و فقط پابرگ را نمایش می‌دهند. این طراحی تمیز،
توجه خواننده را به ابتدای فصل جلب می‌کند.

\section{صفحات معمولی}

صفحات معمولی دارای سربرگ کامل (نام فصل/بخش + شماره صفحه)
و پابرگ هستند. خط افقی سربرگ و پابرگ با ضخامت ۰.۴pt
طراحی شده است.

\section{فاصله خطوط}

فاصله خطوط \lr{MatinBook} روی ۱.۱۵ تنظیم شده است
(\texttt{\textbackslash setstretch\{1.15\}}). این مقدار
برای متن فارسی با فونت \lr{XB Niloofar} بهینه است
و خوانایی را افزایش می‌دهد بدون اینکه فضای زیادی
اشغال کند.

%----------------------------------------
% Chapter 6: Summary
%----------------------------------------
\chapter{خلاصه نتایج}
\label{chap:summary}

\section{وضعیت صفحه‌آرایی}

\begin{table}[h]
\centering
\caption{وضعیت اجزای صفحه‌آرایی}
\label{tab:layout-status}
\begin{tabular}{|c|c|C{6cm}|}
    \hline
    \textbf{بخش} & \textbf{وضعیت} & \textbf{توضیحات} \\
    \hline
    هندسه صفحه & \textcolor{matingreen}{✅ فعال} & A4، حاشیه ۲.۵cm، صحافی ۱cm \\
    سربرگ & \textcolor{matingreen}{✅ فعال} & دوطرفه، با نام فصل و شماره \\
    پابرگ & \textcolor{matingreen}{✅ فعال} & با نام MatinBook \\
    استایل فصل & \textcolor{matingreen}{✅ فعال} & شماره بزرگ + خط رنگی \\
    فهرست مطالب & \textcolor{matingreen}{✅ فعال} & با نقطه‌چین رنگی \\
    پانویس & \textcolor{matingreen}{✅ فعال} & خط ۳۰٪ عرض \\
    فاصله خطوط & \textcolor{matingreen}{✅ فعال} & ۱.۱۵ \\
    \hline
\end{tabular}
\end{table}

\section{نتیجه}

\textbf{ماژول‌های صفحه‌آرایی \lr{MatinBook} نسخه ۱ با موفقیت آزمایش شدند.}
تمام حاشیه‌ها، سربرگ‌ها، پابرگ‌ها و عناوین به درستی کار می‌کنند. 🎉

\end{document}